\section{Related Work}

% ============================================================================
% DRAFT v2 (post-Ali reframe). Register: Laban 2025 / Kamoi 2024.
% 8 themed citation clusters ("sea of blue"). No em-dashes, no staccato.
% Leads with human-AI interaction; over-reliance is the backbone of "hidden cost".
% Each paragraph ties back to our contribution.
% ============================================================================

\paragraph{Human-AI interaction and user-side determinants of quality.}
A growing body of work treats output quality as a property of the interaction rather than
of the model alone \cite{laban2025lost, yang2025underspecification}, and scalable-oversight
research identifies the condition we study: users who cannot articulate precise intent or
validate a complex output \cite{zhou2026oversight}. We locate the revision request itself
as the user-side variable and ask what its most impoverished form does.

\paragraph{Over-reliance and automation bias.}
Whether such degradation is noticed depends on the user's ability to evaluate the output.
Decades of automation research find that people accept automated output they cannot verify
\cite{skitka1999automation, parasuraman2010complacency}, and in AI-assisted decisions
fluent presentation raises acceptance without improving discrimination
\cite{bansal2021whole, bucinca2021trust}. Users rely rather than pay the cost of
verification \cite{vasconcelos2023explanations, schemmer2023reliance}, which is why a quiet
quality decline is absorbed rather than caught.

\paragraph{Intrinsic self-correction and self-refinement.}
Self-refinement work converges on a single finding: intrinsic prompting alone is
unreliable \cite{madaan2024self, huang2024large, kamoi2024selfcorrection, tyen2024llms}.
Reliable self-correction depends instead on a strong external verifier or on additional
training \cite{zhang2024small, qu2024recursive, shinn2023reflexion}. What that literature
asks throughout is whether a given feedback signal suffices to improve an output. We ask
the complementary question of what happens when the signal is absent entirely and the
output to be revised is already adequate.

\paragraph{Sycophancy and compliance under pressure.}
Instruction-tuned models comply and avoid disagreement even when holding firm would help
\cite{perez2023discovering, sharma2024towards, wei2024simple, openai2025sycophancy}, and
under sustained pressure correct answers get revised toward worse ones
\cite{xu2024earth, fanous2025syceval}. Our repeated probe is a weak instance of that
pressure.

\paragraph{Biases in LLM-as-judge evaluation.}
Using LLMs to evaluate generated text is standard practice \cite{zheng2024judging,
li2024generation}, but judges carry self-preference and position biases
\cite{panickssery2024llm, koo2024benchmarking} and reward stylistic and formatting cues
over substance \cite{singhal2023long, wu2023style, zhang2024lists, chen2024humans,
ye2024justice}. Meta-commentary is precisely such a cue, and we find it biases a judge
strongly enough to reverse the measured direction of the effect.

\paragraph{Multi-turn degradation and closest work.}
Quality loss across multiple turns has been observed before through different mechanisms:
early dialogue systems lose coherence through topic drift rather than revision
\cite{zhang2020dialogpt, thoppilan2022lamda}. Closest to our work, \citet{laban2025lost}
show that models lose accuracy when a task's requirements arrive piecewise, and find the
effect threshold-shaped rather than cumulative. Our setting isolates a different cause: the
task is fully specified from the first turn and the initial output is already sufficient,
so the degradation is produced by the revision act itself rather than by information
arriving late. The two findings are complementary. Our revision tax is adjacent to the single-turn over-generation metric of \citet{borisov2026yapbench} but captures a distinct multi-turn cost, the tokens spent revising an already-sufficient output past the point of peak quality.

